\section*{Mock 2 · 20 points}
Let $R=\sqrt{p_1}+\sqrt{p_2}$.

\paragraph{(a), 5 points.}
\[
 v_m=\frac1{R^2},\qquad v_{p_i}=-\frac{m}{\sqrt{p_i}R^3},\qquad
 x_i=-\frac{v_{p_i}}{v_m}=\frac{m}{\sqrt{p_i}R}.
\]
The budget identity is $\sum_i p_ix_i=m(\sqrt{p_1}+\sqrt{p_2})/R=m$.
Suggested split: 2 for derivatives, 2 for demands, 1 for spending.

\paragraph{(b), 5 points.}
Invert $v(p,e)=\bar u$:
\[
 e(p,\bar u)=\bar uR^2,\qquad
 h_i=e_{p_i}=\frac{\bar uR}{\sqrt{p_i}}.
\]
Suggested split: 2 for expenditure, 3 for both Hicksian quantities.

\paragraph{(c), 6 points.}
At old prices, $R^0=2$, $x^0=(4,4)$ and $u^0=2$. At new prices,
$R^1=3$, $x^1=(4/3,8/3)$ and $u^1=8/9$. Therefore
\[
 CV=e((4,1),2)-8=18-8=10,
\]
\[
 EV=8-e((1,1),8/9)=8-32/9=40/9.
\]
Suggested split: 2 for utilities, 2 for each correctly signed monetary measure.

\paragraph{(d), 4 points.}
\[
 L_{CS}=\int_1^4\frac8{\sqrt s(\sqrt s+1)}\,ds
       =16\int_1^2\frac{dt}{t+1}
       =16\log(3/2)\approx6.4874.
\]
Thus $40/9<L_{CS}<10$, or $EV<L_{CS}<CV$. The three amounts differ because
income effects matter. Suggested split: 3 for the integral and value, 1 for comparison.
